\lyseis{Β2.6}{Παράλληλες ευθείες και τέμνουσα}
\sol{1} Τέσσερις γωνίες είναι $55^\circ$ και τέσσερις $125^\circ$.
\sol{2} Εντός εναλλάξ ίσες: $x=2x-30^\circ$, άρα $x=30^\circ$.
\sol{3} Παραπληρωματικές: $x+(x+40^\circ)=180^\circ$, $x=70^\circ$. Οι γωνίες είναι $70^\circ$ και $110^\circ$.
\sol{4} Εκτός εναλλάξ ίσες: $3x=x+50^\circ$, $x=25^\circ$.
\sol{5} α) εντός εναλλάξ. β) εντός επί τα αυτά. γ) εκτός εναλλάξ. δ) εντός--εκτός επί τα αυτά.
\sol{6} Ελεύθερη κατασκευή. Οι γωνίες είναι $72^\circ$ (τέσσερις) και $108^\circ$ (τέσσερις).
\sol{7} α) Λάθος (μόνο για παράλληλες). β) Σωστό. γ) Λάθος (παραπληρωματικές).
\sol{8} $50^\circ$ και $130^\circ$.
\sol{9} $70^\circ$ (εντός--εκτός επί τα αυτά, ίσες).
\sol{10} Τέσσερις γωνίες $110^\circ$ και τέσσερις $70^\circ$.
\sol{11} Ναι, γιατί οι εντός εναλλάξ γωνίες είναι ίσες.
\sol{12} Τα ράφια είναι παράλληλα (και οι δύο κάθετες στον στύλο).
\sol{13} Οι εντός εναλλάξ είναι ίσες, και $\hat{\mathrm{Α}}_3+\hat{\mathrm{Α}}_2=180^\circ$ (διαδοχικές), άρα οι εντός επί τα αυτά είναι παραπληρωματικές.
\sol{14} Οι γωνίες $\hat{\mathrm{Β}}$ και $\hat{\mathrm{Γ}}$ του τριγώνου είναι ίσες με τις εντός εναλλάξ γωνίες στο $\mathrm{Α}$. Μαζί με τη $\hat{\mathrm{Α}}$ σχηματίζουν ευθεία γωνία, άρα $\hat{\mathrm{Α}}+\hat{\mathrm{Β}}+\hat{\mathrm{Γ}}=180^\circ$.
\sol{15} Οι γωνίες του είναι $70^\circ,110^\circ,70^\circ,110^\circ$.
\sol{16} Η διχοτόμος σχηματίζει $28^\circ$ (η μισή των $56^\circ$) με την $\varepsilon_1$.
\sol{17} Ελεύθερη απάντηση.
\sol{18} $\hat{\mathrm{Β}}=\hat{\Delta}=117^\circ$ και $\hat{\Gamma}=63^\circ$.
